\section{Introduction}
\label{sec:introduction}

Every recommender in this portfolio is implicit-feedback matrix
factorization; what differs between them is what gets factorized. The anime
project factorizes a $73\text{k} \times 12\text{k}$ matrix of raw user and
item identities with a BPR loss~\cite{rendle2009bpr}; the MovieLens project
factorizes an explicit rating matrix with truncated SVD. Both share the same
structural limitation: a row or column exists only for an identity observed
during training. A user who signs up the day after the last retraining job,
or a product added to the catalogue this morning, has no row or column at
all, and the only fallback available to an identity-indexed model is
non-personalized popularity.

\subsection{The hybrid alternative}

LightFM~\cite{kula2015} reframes the factorization: instead of learning one
embedding per user and one per item, it learns one embedding per
\emph{feature}, and represents a user or item as the sum of the embeddings
of its known features,
\begin{equation}
  \pvec_u = \sum_{f \in F_u} \evec_f, \qquad
  \qvec_i = \sum_{g \in F_i} \evec_g .
  \label{eq:feature-sum}
\end{equation}
A brand-new user who has filled in a skin-type and skin-tone profile already
has a well-defined $\pvec_u$, built from embeddings the model learned from
every \emph{other} user who shares those two attributes. The same holds for a
brand-new product described by its brand, category, price tier and
ingredient tags. Identity is folded in as just one more feature -- present
for users/items observed during training, deliberately absent otherwise --
so the model degrades gracefully to a pure content-based recommender rather
than failing outright.

This is a classical, linear, SGD-trained model: there are no hidden layers
and no nonlinear activations between the feature sum and the score. The
contribution of this report is not a new loss function; it is (a) a
from-scratch, dependency-free implementation of the WARP loss that LightFM
popularized, and (b) a cold-start evaluation protocol precise enough to
quantify the gap between a hybrid and an identity-only model, rather than
asserting it qualitatively.

\subsection{Why not depend on the \texttt{lightfm} package}

The \texttt{lightfm} PyPI package exists and implements the same idea, but
its Cython extension is compiled against a fixed, old NumPy ABI and routinely
fails to build on current Python/NumPy combinations -- exactly the kind of
fragile dependency a portfolio repository should not ask a reviewer to fight
with. Reimplementing the algorithm in plain NumPy/SciPy (Section
\ref{sec:methodology}) removes that dependency entirely while keeping every
equation in this report directly traceable to a function in the
repository.

\subsection{Scope and claims}

\begin{enumerate}
  \item The feature layer is point-in-time correct. Aggregate features are
        computed with a strict \texttt{<} cutoff against each user's own
        interaction order, never \texttt{<=}, and this is enforced by tests
        that inject a future review and assert no feature value changes
        (Section~\ref{sec:features}).
  \item Cold-start is an operationalized test, not a claim. A population of
        users and items is withheld from training \emph{entirely} -- not
        merely their last interaction -- and the identity-feature vocabulary
        is built only from the training partition, so a cold entity provably
        has no ID token available (Section~\ref{sec:evaluation}).
  \item The comparison is fair. The ID-only baseline is the same family of
        pairwise SGD factorization (Section~\ref{sec:methodology}), not a
        strawman; it loses on the cold-start slice for a structural reason
        (no row/column exists for a withheld entity), not a tuning one.
\end{enumerate}

\subsection{Repository}

The implementation, infrastructure, and the scripts that generate every
table in this document are available in the accompanying repository. The
results tables (Tables~\ref{tab:results} and \ref{tab:cold-start}) are
emitted by \texttt{scripts/export\_results\_table.py} and the feature
catalogue (Table~\ref{tab:features}) by
\texttt{scripts/export\_feature\_tables.py}, so the document cannot drift
away from the code.
